\section{Structural Origin of the Admissibility Hierarchy and the SU(2) Emergence Thread}
  \label{sec:step4}

  \subsection{The ratio $\sqrt{4/3}$ is not accidental}
    \label{sec:step4-origin}

    The admissibility hierarchy established in Proposition~\ref{prop:main},
    \begin{equation}
      \frac{A^{\max}_{\rho_4}}{A^{\max}_{\rho_k}}
      = \sqrt{\frac{8}{6}} = \sqrt{\frac{4}{3}} \approx 1.155,
      \label{eq:ratio}
    \end{equation}
    has a precise internal origin in the $Q_8$ Cayley spectrum.
    It originates entirely from the Laplacian eigenvalues
    $\lambda_{\rho_4} = 6$ and $\lambda_{\rho_k} = 8$.
    The ratio $8/6$ is in turn fixed by the adjacency eigenvalues
    $\mu_{\rho_k} = -2$ and $\mu_{\rho_4} = 0$ through $\lambda = d - \mu$ with $d = 6$.

    The central observation is that $\mu_{\rho_4} = 0$.
    The quaternionic sector lies exactly at the centre of the adjacency spectrum.
    This is a structural consequence of the character theory of $Q_8$,
    not a numerical coincidence.

  \subsection{Why $\mu_{\rho_4} = 0$: character-theoretic argument}
    \label{sec:step4-character}

    For any Cayley graph on a finite group $G$ with symmetric generating set $S$,
    Schur's lemma gives the adjacency eigenvalue in the irreducible sector $\rho$ as
    \begin{equation}
      \mu_\rho = \frac{1}{\dim \rho} \sum_{s \in S} \chi_\rho(s).
      \label{eq:mu-character}
    \end{equation}
    For the quaternionic representation $\rho_4$ of $Q_8$, the generating set
    $S = \{\pm i, \pm j, \pm k\}$ maps to the matrices
    \begin{equation}
      \rho_4({\pm i}) = \pm\begin{pmatrix}
                             i & 0 \\ 0 & -i
      \end{pmatrix},
      \quad
      \rho_4({\pm j}) = \pm\begin{pmatrix}
                             0 & 1 \\ -1 & 0
      \end{pmatrix},
      \quad
      \rho_4({\pm k}) = \pm\begin{pmatrix}
                             0 & i \\ i & 0
      \end{pmatrix}.
    \end{equation}
    All six generator matrices are traceless.
    Therefore
    \begin{equation}
      \chi_{\rho_4}(s) = \mathrm{tr}(\rho_4(s)) = 0
      \qquad
      \forall\, s \in S,
      \label{eq:vanishing-character}
    \end{equation}
    and the character sum in Eq.~\eqref{eq:mu-character} vanishes identically,
    giving $\mu_{\rho_4} = 0$.

    This is a structural property.
    The generators of $Q_8$ map to traceless matrices under $\rho_4$
    because $Q_8$ is generated by the unit quaternions $i,j,k$,
    whose matrix images are exactly the imaginary units of $\mathrm{SU}(2)$.

  \subsection{$Q_8$ as the minimal discrete encoding of spin}
    \label{sec:step4-spin}

    The connection between $Q_8$ and spin is not coincidental but topological.
    The simply-connected double cover of $SO(3)$ is $\mathrm{SU}(2)$,
    which is topologically the three-sphere $S^3$:
    \begin{equation}
      \mathrm{SU}(2)
      \cong
      \big\{a_0 I + i a_1 \sigma_x + i a_2 \sigma_y + i a_3 \sigma_z
      \;:\;
      a_0^2 + a_1^2 + a_2^2 + a_3^2 = 1\big\}
      \cong S^3.
      \label{eq:SU2-S3}
    \end{equation}

    The group $Q_8 = \{\pm 1, \pm i, \pm j, \pm k\}$ is the smallest
    non-abelian subgroup of the unit quaternions, hence the smallest
    non-abelian subgroup of $\mathrm{SU}(2)$.
    Its elements encode the minimal subgroup structure compatible with spinorial
    $\mathrm{SU}(2)$ geometry.

  \subsection{Connection with the $S^3$ spectral invariant $\lambda_2/\lambda_1 = 8/3$}
    \label{sec:step4-S3}

    On the three-sphere $S^3$, the Laplacian eigenvalues are
    $\lambda_\ell = \ell(\ell + 2)$ with multiplicity $(\ell+1)^2$, giving
    \begin{equation}
      \ell = 0:\;\lambda_0 = 0,
      \qquad
      \ell = 1:\;\lambda_1 = 3,
      \qquad
      \ell = 2:\;\lambda_2 = 8,
      \qquad
      \ell = 3:\;\lambda_3 = 15.
    \end{equation}
    The fundamental spectral ratio of $S^3$ is
    \begin{equation}
      \frac{\lambda_2}{\lambda_1}\bigg|_{S^3} = \frac{8}{3}.
      \label{eq:S3-ratio}
    \end{equation}
    This ratio is the primary spectral invariant used elsewhere in the
    Cosmochrony programme to characterise the $S^3$ geometry from the
    relational spectrum~\cite{Spectral121}.

    The ratio from the present $Q_8$ calculation,
    \begin{equation}
      \frac{\lambda_{\rho_k}}{\lambda_{\rho_4}} = \frac{8}{6} = \frac{4}{3},
    \end{equation}
    is related heuristically to the $S^3$ ratio by
    \begin{equation}
      \frac{\lambda_{\rho_k}}{\lambda_{\rho_4}}
      = \frac{1}{2}\cdot\frac{\lambda_2}{\lambda_1}\bigg|_{S^3}.
      \label{eq:ratio-relation}
    \end{equation}

  \subsection{The SU(2) emergence programme}
    \label{sec:step4-programme}

    The $Q_8$ result motivates a systematic programme toward $\mathrm{SU}(2)$
    emergence via spectral admissibility.
    The key chain of inclusions is
    \begin{equation}
      Q_8 \subset 2I \subset \mathrm{SU}(2),
      \label{eq:inclusion-chain}
    \end{equation}
    where $2I$ is the binary icosahedral group of order $120$.

    The asymptotic step is provided by the Ramanujan graphs of
    Lubotzky, Phillips, and Sarnak~\cite{LPS1988}, which are Cayley graphs of
    $\mathrm{SL}(2,\mathbb{F}_p)$ designed to approximate $S^3$ with near-optimal
    spectral gap.

  \subsection{Physical interpretation}
    \label{sec:step4-physics}

    The non-abelian quaternionic sector has zero net correlation with the
    generators of the graph, $\mu_{\rho_4} = 0$.
    This neutrality is the source of its resilience.
    The bounded-flux constraint imposes less penalty on modes that do not couple
    positively or negatively to the relational connectivity.

    The mechanism is therefore a structural consequence of the embedding
    $Q_8 \hookrightarrow \mathrm{SU}(2)$ together with bounded substrate capacity:
    \begin{equation}
      \text{bounded relational capacity}
      \;\Longrightarrow\;
      \text{non-abelian sector resilience}.
    \end{equation}
